% GENERATED FILE. Edit canonical lessons, then run npm run generate:books.
% canonical-source-hash: fnv1a64:85481c3eefaf45e1
% canonical-lessons: JA-C241-sasou, JA-C241-shibai, JA-C241-yukai, JA-C241-moushikomu, JA-C241-yaru, JA-R241-review-236-241

\chapter{Invitations and Plays}
\label{ch:ja-invitations-and-plays}
\hlchaptermodality{241}

\begin{chapteropening}
\textbf{By the end of this chapter:} \emph{I can say who invited me and what we did together.}

Complete the last lesson of chapter 241: I can say who invited me and what we did together.
\end{chapteropening}

\section[\texorpdfstring{\ja{さそう} (sasou)}{さそう (sasou)}]{\ja{さそう} (sasou) — to invite; to ask out}

\label{lesson:JA-C241-sasou}



\begin{quote}
Before the new one: say the Japanese for disliked, then the Japanese for to get tired of.
\end{quote}

\subsection*{You'll want to know: \ja{さそう}}
\textbf{\ja{さそう}} — \emph{sasou} — \textquotedblleft{}to invite; to ask out\textquotedblright{}.

\begin{quote}
\textbf{\ja{せんしゅう} \ja{ともだちを} \ja{えいがに} \ja{さそいました。}} — \emph{senshū tomodachi o eiga ni sasoimashita} — \textquotedblleft{}Last week I invited a friend to a film.\textquotedblright{}
\end{quote}

Polite \textbf{\ja{さそいます}}, \emph{sasoimasu}; a godan verb.

\begin{cousinweb}
\textbf{\ja{さそう}} (\emph{sasou}) is a native Japanese word (wago), inherited rather than borrowed from Chinese; its kanji was attached later as a kun'yomi (JMdict, KANJIDIC2).
\end{cousinweb}

\begin{grammarlens}[title={What you've built}]
One more word for everyday life at A2.
\end{grammarlens}

\subsection*{Your turn}
\begin{itemize}
\raggedright
  \item \emph{Say it:} \emph{sasou}
  \item \emph{Say it:} \emph{sasou}, once more
  \item \emph{Say it:} \emph{senshū tomodachi o eiga ni sasoimashita} — Last week I invited a friend to a film.
  \item \emph{Recall:} say the Japanese for disliked, then the Japanese for to get tired of, then say \emph{sasou} again
\end{itemize}

\begin{tcolorbox}[breakable,colback=teal!4,colframe=teal!35!black,title={Before you move on}]
What does \ja{さそう} mean? (\textquotedblleft{}To invite; to ask out\textquotedblright{}.) Say it once more.
\end{tcolorbox}
\section[\texorpdfstring{\ja{しばい} (shibai)}{しばい (shibai)}]{\ja{しばい} (shibai) — a stage play; theatre}

\label{lesson:JA-C241-shibai}



\begin{quote}
Before the new one: say the Japanese for to get tired of, then the Japanese for to invite.
\end{quote}

\subsection*{You'll want to know: \ja{しばい}}
\textbf{\ja{しばい}} — \emph{shibai} — \textquotedblleft{}a stage play; theatre\textquotedblright{}.

\begin{quote}
\textbf{\ja{せんげつ} \ja{ははと} \ja{しばいを} \ja{みました。}} — \emph{sengetsu haha to shibai o mimashita} — \textquotedblleft{}Last month I saw a play with my mother.\textquotedblright{}
\end{quote}

\begin{cousinweb}
\textbf{\ja{しばい}} (\emph{shibai}) is a native Japanese compound (wago), inherited rather than borrowed: it is built from \emph{shiba}, \textquotedblleft{}turf, lawn\textquotedblright{}, and \emph{i}, from \emph{iru}, \textquotedblleft{}to be, to sit\textquotedblright{} (JMdict, KANJIDIC2).
\end{cousinweb}

\begin{grammarlens}[title={What you've built}]
One more word for everyday life at A2.
\end{grammarlens}

\subsection*{Your turn}
\begin{itemize}
\raggedright
  \item \emph{Say it:} \emph{shibai}
  \item \emph{Say it:} \emph{shibai}, once more
  \item \emph{Say it:} \emph{sengetsu haha to shibai o mimashita} — Last month I saw a play with my mother.
  \item \emph{Recall:} say the Japanese for to get tired of, then the Japanese for to invite, then say \emph{shibai} again
\end{itemize}

\begin{tcolorbox}[breakable,colback=teal!4,colframe=teal!35!black,title={Before you move on}]
What does \ja{しばい} mean? (\textquotedblleft{}A stage play; theatre\textquotedblright{}.) Say it once more.
\end{tcolorbox}
\section[\texorpdfstring{\ja{ゆかい} (yukai)}{ゆかい (yukai)}]{\ja{ゆかい} (yukai) — pleasant; delightful}

\label{lesson:JA-C241-yukai}



\begin{quote}
Before the new one: say the Japanese for to invite, then the Japanese for a stage play.
\end{quote}

\subsection*{You'll want to know: \ja{ゆかい}}
\textbf{\ja{ゆかい}} — \emph{yukai} — \textquotedblleft{}pleasant; delightful\textquotedblright{}.

\begin{quote}
\textbf{\ja{きのうの} \ja{しばいは} \ja{ゆかいでした。}} — \emph{kinō no shibai wa yukai deshita} — \textquotedblleft{}Yesterday's play was delightful.\textquotedblright{}
\end{quote}

\begin{cousinweb}
\textbf{\ja{ゆかい}} (\emph{yukai}) is Sino-Japanese (kango): two Chinese characters read with their on'yomi, readings borrowed from Chinese. They are \emph{yu}, \textquotedblleft{}joy\textquotedblright{} (Middle Chinese \emph{yu}) and \emph{kai}, \textquotedblleft{}cheerful\textquotedblright{} (\emph{khwaejH}) (KANJIDIC2; Middle Chinese in Baxter's transcription).
\end{cousinweb}

\begin{grammarlens}[title={What you've built}]
One more word for everyday life at A2.
\end{grammarlens}

\subsection*{Your turn}
\begin{itemize}
\raggedright
  \item \emph{Say it:} \emph{yukai}
  \item \emph{Say it:} \emph{yukai}, once more
  \item \emph{Say it:} \emph{kinō no shibai wa yukai deshita} — Yesterday's play was delightful.
  \item \emph{Recall:} say the Japanese for to invite, then the Japanese for a stage play, then say \emph{yukai} again
\end{itemize}

\begin{tcolorbox}[breakable,colback=teal!4,colframe=teal!35!black,title={Before you move on}]
What does \ja{ゆかい} mean? (\textquotedblleft{}Pleasant; delightful\textquotedblright{}.) Say it once more.
\end{tcolorbox}
\section[\texorpdfstring{\ja{もうしこむ} (mōshikomu)}{もうしこむ (mōshikomu)}]{\ja{もうしこむ} (mōshikomu) — to sign up; to apply}

\label{lesson:JA-C241-moushikomu}



\begin{quote}
Before the new one: say the Japanese for a stage play, then the Japanese for pleasant.
\end{quote}

\subsection*{You'll want to know: \ja{もうしこむ}}
\textbf{\ja{もうしこむ}} — \emph{mōshikomu} — \textquotedblleft{}to sign up; to apply\textquotedblright{}.

\begin{quote}
\textbf{\ja{せんげつ} \ja{あには} \ja{じゅうどうの} \ja{しあいに} \ja{もうしこみました。}} — \emph{sengetsu ani wa jūdō no shiai ni mōshikomimashita} — \textquotedblleft{}Last month my elder brother signed up for a judo match.\textquotedblright{}
\end{quote}

Polite \textbf{\ja{もうしこみます}}, \emph{mōshikomimasu}; a godan verb.

\begin{cousinweb}
\textbf{\ja{もうしこむ}} (\emph{mōshikomu}) is a native Japanese compound (wago), inherited rather than borrowed: it is built from \emph{mōshi}, from \emph{mōsu}, \textquotedblleft{}to say\textquotedblright{}, and \emph{komu}, \textquotedblleft{}to go in, to crowd in\textquotedblright{} (JMdict, KANJIDIC2).
\end{cousinweb}

\begin{grammarlens}[title={What you've built}]
One more word for everyday life at A2.
\end{grammarlens}

\subsection*{Your turn}
\begin{itemize}
\raggedright
  \item \emph{Say it:} \emph{mōshikomu}
  \item \emph{Say it:} \emph{mōshikomu}, once more
  \item \emph{Say it:} \emph{sengetsu ani wa jūdō no shiai ni mōshikomimashita} — Last month my elder brother signed up for a judo match.
  \item \emph{Recall:} say the Japanese for a stage play, then the Japanese for pleasant, then say \emph{mōshikomu} again
\end{itemize}

\begin{tcolorbox}[breakable,colback=teal!4,colframe=teal!35!black,title={Before you move on}]
What does \ja{もうしこむ} mean? (\textquotedblleft{}To sign up; to apply\textquotedblright{}.) Say it once more.
\end{tcolorbox}
\section[\texorpdfstring{\ja{やる} (yaru)}{やる (yaru)}]{\ja{やる} (yaru) — to do; to play (a game)}

\label{lesson:JA-C241-yaru}



\begin{quote}
Before the new one: say the Japanese for pleasant, then the Japanese for to sign up.
\end{quote}

\subsection*{You'll want to know: \ja{やる}}
\textbf{\ja{やる}} — \emph{yaru} — \textquotedblleft{}to do; to play (a game)\textquotedblright{}.

\begin{quote}
\textbf{\ja{きのう} \ja{ともだちと} \ja{しょうぎを} \ja{やりました。}} — \emph{kinō tomodachi to shōgi o yarimashita} — \textquotedblleft{}Yesterday I played shogi with a friend.\textquotedblright{}
\end{quote}

Polite \textbf{\ja{やります}}, \emph{yarimasu}; a godan verb, more casual than \textbf{\ja{します}}.

\begin{cousinweb}
\textbf{\ja{やる}} (\emph{yaru}) is a native Japanese word (wago), inherited rather than borrowed from Chinese; its kanji was attached later as a kun'yomi, and the word is usually written in kana (JMdict, KANJIDIC2).
\end{cousinweb}

\begin{grammarlens}[title={What you've built}]
One more word for everyday life at A2.
\end{grammarlens}

\subsection*{Your turn}
\begin{itemize}
\raggedright
  \item \emph{Say it:} \emph{yaru}
  \item \emph{Say it:} \emph{yaru}, once more
  \item \emph{Say it:} \emph{kinō tomodachi to shōgi o yarimashita} — Yesterday I played shogi with a friend.
  \item \emph{Recall:} say the Japanese for pleasant, then the Japanese for to sign up, then say \emph{yaru} again
\end{itemize}

\begin{tcolorbox}[breakable,colback=teal!4,colframe=teal!35!black,title={Before you move on}]
What does \ja{やる} mean? (\textquotedblleft{}To do; to play (a game)\textquotedblright{}.) Say it once more.
\end{tcolorbox}
\section[(dialogue)]{Words from chapters 236–241 — a review}

\label{lesson:JA-R241-review-236-241}



\begin{quote}
Say the Japanese for a film and to watch.
\end{quote}

\subsection*{Your turn}
Say these aloud:

\begin{itemize}
\raggedright
  \item a film, to watch, to listen, music, a musical performance
  \item a photo, to photograph, to collect, to gather, a doll
  \item comics, a novel, reading, a programme, boring
  \item to dance, to go fishing, a walk, a magic trick, shogi
  \item skilful, unskilful, liked, disliked, to get tired of
  \item to invite, a stage play, pleasant, to sign up, to do
\end{itemize}

\begin{tcolorbox}[breakable,colback=teal!4,colframe=teal!35!black,title={Before you move on}]
A film? (\textbf{\ja{えいが}}, \emph{eiga}.) To watch? (\textbf{\ja{みる}}, \emph{miru}.) To listen? (\textbf{\ja{きく}}, \emph{kiku}.) Music? (\textbf{\ja{おんがく}}, \emph{ongaku}.) A musical performance? (\textbf{\ja{えんそう}}, \emph{ensō}.) A photo? (\textbf{\ja{しゃしん}}, \emph{shashin}.) To photograph? (\textbf{\ja{とる}}, \emph{toru}.) To collect? (\textbf{\ja{あつめる}}, \emph{atsumeru}.) To gather? (\textbf{\ja{あつまる}}, \emph{atsumaru}.) A doll? (\textbf{\ja{にんぎょう}}, \emph{ningyō}.) Comics? (\textbf{\ja{まんが}}, \emph{manga}.) A novel? (\textbf{\ja{しょうせつ}}, \emph{shōsetsu}.) Reading? (\textbf{\ja{どくしょ}}, \emph{dokusho}.) A programme? (\textbf{\ja{ばんぐみ}}, \emph{bangumi}.) Boring? (\textbf{\ja{つまらない}}, \emph{tsumaranai}.) To dance? (\textbf{\ja{おどる}}, \emph{odoru}.) To go fishing? (\textbf{\ja{つる}}, \emph{tsuru}.) A walk? (\textbf{\ja{さんぽ}}, \emph{sanpo}.) A magic trick? (\textbf{\ja{てじな}}, \emph{tejina}.) Shogi? (\textbf{\ja{しょうぎ}}, \emph{shōgi}.) Skilful? (\textbf{\ja{じょうず}}, \emph{jōzu}.) Unskilful? (\textbf{\ja{へた}}, \emph{heta}.) Liked? (\textbf{\ja{すき}}, \emph{suki}.) Disliked? (\textbf{\ja{きらい}}, \emph{kirai}.) To get tired of? (\textbf{\ja{あきる}}, \emph{akiru}.) To invite? (\textbf{\ja{さそう}}, \emph{sasou}.) A stage play? (\textbf{\ja{しばい}}, \emph{shibai}.) Pleasant? (\textbf{\ja{ゆかい}}, \emph{yukai}.) To sign up? (\textbf{\ja{もうしこむ}}, \emph{mōshikomu}.) To do? (\textbf{\ja{やる}}, \emph{yaru}.)
\end{tcolorbox}
